% Generated draft: requires mathematical and source review.
\documentclass[11pt]{article}
\usepackage[margin=25mm]{geometry}
\usepackage{fontspec}
\setmainfont{Latin Modern Roman}
\usepackage{amsmath,amssymb,mathtools,unicode-math}
\setmathfont{Latin Modern Math}
\usepackage{enumitem,needspace,xurl,hyperref}
\setlength{\parindent}{0pt}
\setlength{\parskip}{5pt}
\setlength{\emergencystretch}{3em}
\newcommand{\N}{\mathbb{N}}
\newcommand{\Z}{\mathbb{Z}}
\newcommand{\Q}{\mathbb{Q}}
\newcommand{\R}{\mathbb{R}}
\newcommand{\C}{\mathbb{C}}
\newcommand{\F}{\mathbb{F}}
\newcommand{\A}{\mathbb{A}}
\newcommand{\PP}{\mathbb{P}}
\newcommand{\E}{\mathbb{E}}
\newcommand{\eps}{\varepsilon}
\newcommand{\dd}{\,\mathrm{d}}
\begin{document}
\section*{UnsolvedMath ID: 801085}
\textbf{Source problem code:} catalogue entry
\subsection{Definitions and mathematical statement}
Let \(J\) denote the Julia set referred to in the source question, and let \(f_*\) denote the source's operator on functions on \(J\). The supplied record does not specify the map \(f\), the definition of \(J\), or the action and normalization of \(f_*\). For a potential \(\phi\), write formally
\[
\mathcal L_\phi g=f_*\bigl(e^{\phi-P(\phi)}g\bigr),
\]
whenever these operations are defined. The admissible potentials and the quantity \(P(\phi)\) are not defined in the supplied record. An invariant line for an operator \(T\) means a one-dimensional linear subspace \(E\) satisfying \(T(E)\subseteq E\).

\textbf{Question (Berger and Bedford, Question 37, 2023).} Do there exist norms for functions on \(J\) which are bounded on H\"older continuous functions, contract on the complement of an invariant line under \(f_*\), or, more generally, under \(\mathcal L_\phi\), and for which the contraction is stable under perturbation of \(\phi\)?

The question retains unresolved qualifications from the supplied source. The function space, scalar field, metric on \(J\), H\"older exponent, meaning of boundedness, invariant line, and meaning of its complement are unspecified. For example, a continuous-inclusion interpretation of boundedness would require an estimate \(\|g\|\leq C\|g\|_{C^\alpha(J)}\) for specified \(\alpha\), metric, and constant \(C\); this interpretation is not fixed by the record. Likewise, contraction could refer to an operator-norm bound strictly below one on a specified invariant complement, decay of iterates, or contraction of an induced quotient operator. The record does not select among these meanings. The topology of perturbations of \(\phi\), the permitted dependence of the norm and invariant line on \(\phi\), and any uniformity of contraction constants also remain unspecified. Consequently, this is a statement-only transcription of the existence question, requiring these choices before a fully quantified formulation can be given.
\subsection{Short English statement}
Berger and Bedford ask whether suitable norms on functions on the Julia set make the transfer operator, or its displayed potential-weighted version, contract away from an invariant line, with contraction persisting under perturbations of the potential. The supplied record leaves essential operator, function-space, and stability specifications unresolved.
\subsection{Sources}
\begin{itemize}
\item Local UnsolvedMath registry, permanent ID \(801085\), exact source problem code \texttt{catalogue entry}, titled \emph{Henon Maps: Open Problems --- Question 37}. The record attributes the question to Berger and Bedford, with proposed year 2023, and labels it \texttt{open}; this is a supplied status label, not an independently verified mathematical or literature assessment.
\item Berger and Bedford, \emph{H\'enon Maps: A List of Open Problems} (2023), Question 37. Supplied source URL: \url{https://arxiv.org/abs/2312.03907}. The provenance records access on 2026-10-16 and extraction by \texttt{arxiv\_latex}; these are supplied provenance details, not a claim of independent source access or verification.
\end{itemize}
\end{document}
